% !PRINTED 666
\weberpage{666}

%% [running head, not body text -- commented out 2026-08-07] Page furniture (running
%% head, folio, signature) is excluded from this edition by standing policy.
%% 666 \hfill \gesperrt{Zweiundzwanzigster Abschnitt.} \hfill \S.\,183.

\begin{center}
\S.\,183.

\gesperrt{Beweis des zweiten H\"ulfssatzes f\"ur ein ungerades $m$.}
\end{center}

Das zweite Lemma, was nach \S.\,181 noch zu beweisen ist, ist
folgendes:

2. \gesperrt{Ist $\varepsilon_n$ ein System conjugirter numerischer
Einheiten des K\"orpers $\Omega_m$, das der Bedingung gen\"ugt, dass}

\weq{1}{\varepsilon_n \varepsilon_1^{-n} = \mathfrak{E}^m}

\gesperrt{die $m^{\text{te}}$ Potenz einer Einheit in $\Omega_m$ ist,
so sind die $\varepsilon_n$ selbst $m^{\text{te}}$ Potenzen von
Einheiten in $\Omega_m$, multiplicirt mit einer Einheitswurzel der
Form $r^{nk}$.}

Auch dieser Beweis ist f\"ur den Fall eines ungeraden $m$ ganz
elementar. Wir betrachten daher hier zun\"achst diesen Fall.

Nach \S.\,178, 1. k\"onnen wir jede Einheit des K\"orpers $\Omega_m$
in der Form darstellen:

\weq{2}{\varepsilon = \pm r^k \mathfrak{E}(r), \quad \varepsilon_n = \pm r^{nk} \mathfrak{E}(r^n),}

worin $\mathfrak{E}(r)$ eine reelle Einheit des K\"orpers $\Omega_m$
ist. Es gen\"ugt also $\mathfrak{E}(r)$ der Bedingung

\weq{3}{\mathfrak{E}(r) = \mathfrak{E}(r^{-1}),}

woraus sich nach (2) ergiebt:

\weq{4}{\varepsilon_1 \varepsilon_{-1} = \mathfrak{E}(r)^2;}

aus (1) aber findet sich, wenn man $n = -1$ setzt,

\[
\varepsilon_1 \varepsilon_{-1} = \mathfrak{E}^m,
\]

d.\,h. die linke Seite von (4) ist die $m^{\text{te}}$ Potenz einer
Einheit in $\Omega_m$, und es ist also auch

\weq{5}{\mathfrak{E}(r)^2 = \mathfrak{E}^m}

die $m^{\text{te}}$ Potenz einer solchen Einheit.

Da nun $m$ ungerade ist, so l\"asst sich die ganze rationale Zahl $x$
so bestimmen, dass

\weq{6}{2 x + m = 1}

wird, und daraus folgt

\[
\mathfrak{E}(r) = \mathfrak{E}(r)^{2x} \mathfrak{E}(r)^m.
\]
%CONT>>
